\section{Congruence fibers over discrete tori}\label{sec:geometry}

We construct a graph with two torus ends and a long central region of constant
fiber size. The fiber quotient changes one congruence level at a time outside
this region. These gradual changes keep the degrees bounded, while the long
central region will isolate the topological obstruction at any fixed Rips
scale.

\subsection{The finite quotients and their twists}

Fix a finite symmetric generating set $S$ of
$\Gamma=\SL_3(\Z)$. We use Kazhdan's theorem that $\Gamma$ has property~(T)
\cite{Kazhdan67,BHV08}. For $k\geq 1$, let $Q_k$ be the image of $\Gamma$
under entrywise reduction modulo $2^k$, and put $Q_0=\{1\}$. The natural
reduction maps will always be denoted by their source and target.

\begin{lemma}\label{geom:quotients}
The maps $Q_{k+1}\to Q_k$ are surjective and have kernels of size at most
$2^9$. For every fixed word radius $R$, the reduction $\Gamma\to Q_k$ is
injective on the $S$-word ball of radius $R$ for all sufficiently large $k$.
\end{lemma}

\begin{proof}
Surjectivity follows from the definition as images of the same group. For
$k\geq1$, each matrix modulo $2^k$ has at most $2^9$ entrywise lifts modulo
$2^{k+1}$. For $k=0$, the same bound follows by counting all binary
$3\times3$ matrices. Finally, a word ball is a finite set of integral
matrices. If two of its elements differ, some difference of entries is a
nonzero integer and is not divisible by arbitrarily large powers of two.
Taking the maximum of the finitely many required levels proves injectivity.
\end{proof}

For each $j\geq1$, set
\begin{equation}\label{geom:parameters}
 m_j=2^j,\qquad l_j=100(j+1),\qquad
 U_j=1+m_jE_{12},\qquad V_j=1+m_jE_{23},
\end{equation}
where $E_{ab}$ denotes a matrix unit. A direct multiplication gives
\begin{equation}\label{geom:commutator}
 U_jV_jU_j^{-1}V_j^{-1}=1+m_j^2E_{13}.
\end{equation}
Thus $U_j$ and $V_j$ commute in $Q_{2j}$ and in every lower quotient, although
their commutator remains visible in $Q_{3j}$. This difference between the two
levels is the source of the divisibility argument in
Section~\ref{sec:persistence}.

Let $B_j=(\Z/l_j\Z)^3$ with its axis-neighbor graph. Represent each
coordinate by an integer in $\{0,\ldots,l_j-1\}$. An oriented edge has
signed seam numbers $(a,b,c)\in\Z^3$: a positive step from $l_j-1$ to $0$
contributes $1$ in that coordinate, a reverse step contributes $-1$, and
all other contributions vanish. Equivalently, the unit displacement along
the edge is the difference of the chosen representatives plus
$l_j(a,b,c)$. For $0\leq k\leq2j$, define
\begin{equation}\label{geom:transport}
 g_{j,k}(a,b)=V_j^bU_j^a\pmod {2^k}\in Q_k,
\end{equation}
with the evident interpretation in $Q_0$. Only the first two circle
coordinates will twist the fibers.

The graph $Y_{j,k}$ has vertex set $B_j\times Q_k$. Its edges are of two
types:
\begin{itemize}
\item over a fixed $u\in B_j$, connect $(u,q)$ to $(u,sq)$ for $s\in S$;
\item for an oriented base edge $u\to v$ with seam numbers $(a,b,c)$,
connect $(u,q)$ to $(v,qg_{j,k}(a,b))$.
\end{itemize}
The first action is on the left and the second on the right, so they commute.
Reversing a base edge reverses its matching, because $U_j$ and $V_j$ commute
in $Q_k$. We discard loops and repeated edges. Each fiber Cayley graph is
connected, and the horizontal matchings connect the fibers over the
connected graph $B_j$. Hence $Y_{j,k}$ is connected. Write
\[
 Y_j=Y_{j,2j}.
\]
The metric on $Y_j$ is its intrinsic graph metric.

\subsection{The two ends and the central collar}

For each integer $t$ with $-3j\leq t\leq3j$, put
\begin{equation}\label{geom:height-profile}
 k_j(t)=
 \begin{cases}
  2j,& |t|\leq j,\\
  3j-|t|,& j\leq |t|\leq3j.
 \end{cases}
\end{equation}
At height $t$, place a copy of $Y_{j,k_j(t)}$. Connect successive slices over
the same base vertex by the reduction map from the larger quotient to the
smaller one, or by the identity when the levels agree. The resulting finite
graph is $X_j$.

At height $\pm3j$ the quotient is trivial, so each end is a copy of $B_j$.
The induced graph on the heights $-j,\ldots,j$ is the product of $Y_j$ with
the interval graph: its horizontal edges are those of $Y_j$, and its
vertical edges preserve both $u$ and $q$. We call this subgraph the
\emph{central collar}. Paths in $X_j$ may take shortcuts outside the
collar; whenever we use its product structure, we will first show that the
paths under consideration cannot leave it.

\begin{figure}[htbp]
\centering
\begin{tikzpicture}[x=1.45cm,y=1.05cm,>=Stealth,font=\small]
 \fill[blue!7] (-1,-0.08) rectangle (1,2.27);
 \draw[->,gray] (-3.45,0)--(3.6,0) node[right] {$t$};
 \draw[very thick,blue!65!black] (-3,0)--(-1,2)--(1,2)--(3,0);
 \draw[dashed,gray] (-1,0)--(-1,2);
 \draw[dashed,gray] (1,0)--(1,2);
 \draw[dashed,gray] (0,0)--(0,2);
 \foreach \x/\lab in {-3/{-3j},-1/{-j},0/{0},1/{j},3/{3j}}{
  \draw (\x,0.05)--(\x,-0.05);
  \node[below] at (\x,-0.06) {$\lab$};
 }
 \fill (-3,0) circle (1.7pt);
 \fill (3,0) circle (1.7pt);
 \fill (0,2) circle (1.7pt);
 \node[above left] at (-3,0.08) {$B_j$};
 \node[above right] at (3,0.08) {$B_j$};
 \node[above] at (0,2.07) {$Y_j$ (fiber $Q_{2j}$)};
 \node at (0,1.08) {central collar};
 \node[rotate=45,above] at (-2,1) {$Q_{k_j(t)}$};
 \node[rotate=-45,above] at (2,1) {$Q_{k_j(t)}$};
 \node[anchor=east] at (-3.35,2) {$k_j(t)=2j$};
\end{tikzpicture}
\caption{The congruence level as a function of height in $X_j$.
Every slice lies over the same discrete torus $B_j$. Outside the central
collar, each vertical step changes the quotient level by one. The drawing
records the level profile, not a geometric embedding of the graph or its
fiber cardinalities.}
\label{fig:profile}
\end{figure}

The gradual reductions in Figure~\ref{fig:profile} have bounded fibers by
Lemma~\ref{geom:quotients}. They therefore preserve a uniform degree bound
even though the middle quotient grows rapidly with $j$.

\begin{lemma}\label{geom:basic}
The graphs $B_j,Y_j,X_j$ are finite and connected, with vertex degrees bounded
independently of $j$. Their projections to $B_j$ are $1$-Lipschitz. Each end
copy of $B_j$ is isometrically embedded in $X_j$.
\end{lemma}

\begin{proof}
Only the claims about $X_j$ remain to be checked. Consecutive slices are
connected by surjective reduction maps, so $X_j$ is connected. A vertex has
at most $|S|$ fiber neighbors, six base neighbors, and at most $2^9$ vertical
neighbors on either side. Thus $|S|+6+2^{10}$ is a common bound. Projection
sends an edge to a base edge or a point, proving the Lipschitz assertion.
Within an end, every base path is available. Projection gives the opposite
distance inequality, so the end metric is exactly the metric of $B_j$.
\end{proof}

\subsection{Coarse disjoint unions}

We now fix metrics on
\[
 B=\bigsqcup_{j\geq1}B_j,\qquad
 Y=\bigsqcup_{j\geq1}Y_j,\qquad
 X=\bigsqcup_{j\geq1}X_j.
\]
In $B_j$ choose the zero vertex, in $Y_j$ the vertex $(0,1)$, and in $X_j$
the vertex $(0,1)$ at height zero. In each union, join the chosen vertices
of components $j$ and $j+1$ by an auxiliary edge of length $100+j^2$, and
use the resulting path distances on the original vertex set. These edges
define the metric; they add no vertices to the space.

A path which leaves and returns to the same component must pass through
its chosen vertex both times. Consequently its use of the auxiliary edges
cannot shorten an intrinsic component distance. At each fixed distance
bound, cross-component pairs occur in only finitely many finite components.
The resulting metrics are uniformly discrete and proper. They have bounded
geometry because the weighted graphs have uniformly bounded degrees and
all edge lengths are at least one: a ball of radius $R$ contains no more
vertices than a ball of graph radius $\lfloor R\rfloor$ in a graph of the
same degree bound.

Projection and end inclusion define maps
\begin{equation}\label{geom:coarse-maps}
 f:Y\longrightarrow B,\qquad i_\pm:B\longrightarrow X.
\end{equation}
On each component they are $1$-Lipschitz. For a fixed input distance bound,
the remaining pairs lie in a finite set, so their image distances also have
a finite bound. Thus the maps are bornologous. Preimages of bounded sets
are finite, so the maps are proper. At scale $3$ none of the Rips complexes
has a cross-component simplex. In particular, the end maps induce proper
simplicial inclusions
\[
 i_{\pm,3}:P_3(B)\longrightarrow P_3(X).
\]
These are the maps used to form the K-homology class. The next section shows
why its two endpoint contributions remain distinct at every fixed scale.
